\subsection{Worked example: attitude from vector observations by Gauss--Newton on SO(3)}\label{sec:lie:wahba}

An IMU or star tracker observes known world directions in its own frame: gravity from the accelerometer, magnetic north from a magnetometer, star or landmark bearings from a camera. With world vectors $a_i$ and body-frame measurements $y_i=R\T a_i+n_i$, find $R$:
\[
\min_{R\in\SO}\ \sum_i\|y_i-R\T a_i\|^2 .
\]
\emph{Model.} Residual $r_i(R)=y_i-R\T a_i$. By the inverse-action block of \S\ref{sec:lie:elementary}, $\partial(R\T a_i)/\partial R=\skw{R\T a_i}$, so the residual Jacobian is $-\skw{R\T a_i}$. \emph{Gauss--Newton step}: stack the residuals, solve $(J\T J)\delta=-J\T r$ for $\delta\in\RR^3$, and update $R\leftarrow R\oplus\delta$. The state never leaves $\SO$, and no constraint has to be enforced.

\begin{example}{four directions, 0.01 noise, started 107$^\circ$ away}
Four unit directions (gravity, a tilted magnetic field, two landmark bearings), true attitude $\Exp(0.4,-0.9,2.2)$, noise $\sigma=0.01$, initial guess $R=I$:
\begin{center}\small
\begin{tabular}{@{}cccc@{}}\toprule
iteration & error to truth & $\|\delta\|$ & cost before step\\\midrule
1 & $106.66^\circ$ & 0.975 & 8.15\\
2 & $53.81^\circ$ & 0.993 & 3.74\\
3 & $8.11^\circ$ & 0.801 & 1.09\\
4 & $0.384^\circ$ & 0.136 & $2.7\times10^{-2}$\\
5 & $0.302^\circ$ & $1.8\times10^{-3}$ & $1.04\times10^{-3}$\\
6 & $0.301^\circ$ & $9.6\times10^{-6}$ & $1.04\times10^{-3}$\\\bottomrule
\end{tabular}
\end{center}
Five iterations reach the noise floor ($0.30^\circ$), with the quadratic convergence typical of Gauss--Newton near the optimum. The result agrees with the closed-form SVD (Wahba/Kabsch) solution to $3\times10^{-10}$ degrees. The SVD trick exists only for this particular cost. The Lie-group Gauss--Newton recipe (residual, chain-rule Jacobian, $\oplus$ update) works unchanged for any cost, which is why SLAM back-ends are built on it.
\end{example}
